\chapter{Báo cáo thực nghiệm bổ sung cho SCRFD}

\section{Mục tiêu và thiết lập thực nghiệm}

Phần này chuyển nội dung phân tích thực nghiệm từ mã nguồn và báo cáo Markdown sang định dạng LaTeX của báo cáo chính.
Mục tiêu của cụm thí nghiệm là trả lời hai câu hỏi:
\begin{itemize}
  \item vì sao cải tiến \textbf{ASR+JSAR} tốt hơn \textbf{baseline} trong cùng kiến trúc SCRFD 2.5G;
  \item vì sao khi thay static scale pool mặc định bằng scale set \textbf{SR12} lấy từ phần search của bài báo, mô hình lại thay đổi hành vi theo cách khác.
\end{itemize}

Hai cặp thí nghiệm được so sánh là:
\begin{itemize}
  \item \textbf{Default source scale set}: dùng static scale pool mặc định trong mã nguồn công khai của SCRFD.
  \item \textbf{Paper SR12 scale set}: thay scale pool tĩnh bằng bộ 12 mức được mô tả ở phần sample redistribution của bài báo SCRFD. \cite{guo2022scrfd}
\end{itemize}

Trong cả hai cặp, phần Backbone--Neck--Head giữ nguyên, mô hình được huấn luyện trong \(80\) epochs với cosine scheduler.
Biến độc lập chính chỉ là:
\begin{itemize}
  \item static scale pool ban đầu;
  \item việc bật hay tắt hai cải tiến \textbf{Adaptive Sample Redistribution (ASR)} và \textbf{Joint SampleAssignment Redistribution (JSAR)}.
\end{itemize}

\section{Phương pháp cải tiến}

\subsection{Adaptive Sample Redistribution (ASR)}

ASR thay việc chọn scale augmentation tĩnh bằng một phân phối xác suất động, được cập nhật trực tiếp từ thống kê của quá trình huấn luyện.
Ý tưởng trực quan của ASR được minh hoạ trong Hình~\ref{fig:final-term-asr-schematic}.

\begin{figure}[H]
  \centering
  \includegraphics[width=\textwidth]{final_term_spos_scrfd/asr_schematic.png}
  \caption{Minh hoạ ý tưởng của ASR. Trái: static scale pool ban đầu và ý nghĩa hình học của scale. Giữa: vòng phản hồi từ thống kê huấn luyện sang difficulty của từng size bin, rồi quay lại cập nhật scale policy. Phải: phân phối kích thước mặt sau phép crop + resize dịch về phía nhiều tiny/small samples hơn.}
  \label{fig:final-term-asr-schematic}
\end{figure}

Trong implementation hiện tại, sau mỗi epoch hệ thống gom thống kê theo size bin \(b\):
\begin{itemize}
  \item \(G_b\): số ground-truth faces thuộc bin \(b\);
  \item \(P_b\): số positive anchors khớp với bin \(b\);
  \item \(L^{cls}_b\): tổng classification loss;
  \item \(L^{box}_b\): tổng box loss.
\end{itemize}

Từ đó, ASR tạo ba thành phần difficulty:
\begin{equation}
r_b = \max\!\left(0,\; 1 - \min\!\left(\frac{P_b}{\max(G_b,1)},\; 1\right)\right)
\end{equation}
\begin{equation}
c_b = \frac{L^{cls}_b}{\max(P_b,1)}, \qquad
u_b = \frac{L^{box}_b}{\max(P_b,1)}
\end{equation}

Trong đó \(r_b\) đóng vai trò recall gap proxy, còn \(c_b\) và \(u_b\) đo loss trung bình trên mỗi positive sample.
Với chế độ \texttt{loss\_recall}, difficulty của từng bin được tính bằng tổng của ba đại lượng đã chuẩn hoá:
\begin{equation}
\operatorname{norm}(x)_b = \frac{x_b}{\operatorname{mean}(x_{x>0}) + \varepsilon}
\end{equation}
\begin{equation}
d_b \propto \operatorname{norm}(r_b) + \operatorname{norm}(c_b) + \operatorname{norm}(u_b)
\end{equation}

Difficulty theo bin sau đó được chiếu sang mỗi scale candidate \(s\) thông qua một support prior:
\begin{equation}
S(s,b) = \exp\!\left(-\frac{|\log s - \log \mu_b|}{0.45}\right)
\end{equation}
với \(\mu_b\) là tâm ưu tiên của từng size bin.
Xác suất thô của mỗi scale được xác định bởi:
\begin{equation}
q(s) \propto \sum_b S(s,b)\, d_b
\end{equation}

Sau bước probability floor và EMA smoothing:
\begin{equation}
q'(s) = q(s)(1-Km) + m
\end{equation}
\begin{equation}
p_{t+1}(s) = \alpha\, p_t(s) + (1-\alpha)\, q'(s)
\end{equation}
trong đó \(K\) là số scale candidates, \(m\) là xác suất sàn tối thiểu, còn \(\alpha\) là hệ số làm mượt.

Điểm quan trọng của pipeline SCRFD là:
\begin{itemize}
  \item \textbf{scale nhỏ} \(\Rightarrow\) crop nhỏ hơn \(\Rightarrow\) zoom-in \(\Rightarrow\) mặt lớn hơn sau resize;
  \item \textbf{scale lớn} \(\Rightarrow\) crop rộng hơn \(\Rightarrow\) zoom-out \(\Rightarrow\) mặt nhỏ hơn sau resize.
\end{itemize}

Do đó, khi ASR tăng xác suất ở các scale lớn, phân phối mẫu sau augmentation sẽ dịch về phía nhiều tiny faces hơn.
Nói cách khác, ASR không chỉ đổi cách lấy mẫu trên trục scale, mà đổi cả phân phối khó/dễ của tập huấn luyện mà detector thường xuyên nhìn thấy.

\subsection{Joint SampleAssignment Redistribution (JSAR)}

JSAR không thay đổi kiến trúc suy luận, mà chỉ can thiệp vào bước \textit{assignment} trong lúc huấn luyện để tiny/small faces nhận được supervision dày hơn.
Trực quan của phương pháp được mô tả ở Hình~\ref{fig:final-term-jsar-schematic}.

\begin{figure}[H]
  \centering
  \includegraphics[width=\textwidth]{final_term_spos_scrfd/jsar_schematic.png}
  \caption{Minh hoạ ý tưởng của JSAR. Trái: với tiny GT, ATSS gốc thường chỉ tạo rất ít positives. Giữa: threshold và center gating được nới theo kích thước để nhiều anchors gần tiny GT còn hợp lệ hơn. Phải: nếu số positives vẫn thấp, cơ chế fallback sẽ lấy thêm top-\(k\) anchors nền gần GT nhất theo matching score.}
  \label{fig:final-term-jsar-schematic}
\end{figure}

JSAR được cài trên \texttt{ATSSAssigner} theo ba bước.

\paragraph{Bước 1: nới threshold theo kích thước GT}
ATSS gốc dùng ngưỡng:
\begin{equation}
\tau_g = \mu_g + \sigma_g
\end{equation}
JSAR thay bằng:
\begin{equation}
\tau'_g = \tau_g - \Delta_g
\end{equation}
trong đó:
\begin{equation}
\Delta_g =
\begin{cases}
\delta_{tiny}, & \text{nếu } \operatorname{size}_g < T_{tiny} \\
\delta_{small}, & \text{nếu } T_{tiny} \le \operatorname{size}_g < T_{small} \\
0, & \text{nếu } \operatorname{size}_g \ge T_{small}
\end{cases}
\end{equation}

\paragraph{Bước 2: center gating theo kích thước}
Center gating được làm size-aware bằng hệ số:
\begin{equation}
c_g = \operatorname{clip}\!\left(\frac{\operatorname{size}_g}{T_{tiny}}, 1, \rho\right)
\end{equation}
\begin{equation}
\operatorname{valid}(a,g):\ d_{min}(a,g)\, c_g > t_{center}
\end{equation}
Thiết kế này giúp bớt loại bỏ quá sớm các anchors gần tiny GT.

\paragraph{Bước 3: hybrid fallback khi GT vẫn thiếu positives}
Nếu sau hai bước trên mà tiny/small GT vẫn có quá ít positives:
\begin{equation}
N_{pos}(g) < N_{min}
\end{equation}
thì JSAR lấy thêm anchors nền lân cận theo score:
\begin{equation}
\operatorname{score}(a,g) =
\operatorname{IoU}(a,g) -
0.05 \cdot \frac{\operatorname{dist}(a,g)}{\max(\operatorname{size}_g,1)}
\end{equation}

Trong biến thể soft weighting, trọng số mềm còn có dạng:
\begin{equation}
w(a,g) =
\exp\!\left(-\frac{\operatorname{dist}_{norm}}{T}\right)
\cdot
\operatorname{clamp}\!\left(\operatorname{IoU}(a,g)+0.1,\ 0.05,\ 1.0\right)
\end{equation}

Tuy nhiên trong các run hiện tại, cơ chế chủ đạo là \textit{hybrid fallback}:
\begin{itemize}
  \item nới threshold cho tiny/small GT;
  \item nới center gating theo kích thước;
  \item nếu GT vẫn thiếu positives thì thêm top-\(k\) anchors nền gần nhất theo score.
\end{itemize}

\section{Kết quả định lượng}

\subsection{Metric tổng hợp}

Hình~\ref{fig:final-term-metrics-by-scale-set} và Bảng~\ref{tab:final-term-scrfd-metrics} tóm tắt kết quả WIDERFace của hai scale set.

\begin{figure}[H]
  \centering
  \includegraphics[width=0.88\textwidth]{final_term_spos_scrfd/metrics_by_scale_set.png}
  \caption{So sánh các chỉ số Easy AP, Medium AP, Hard AP và mAP giữa baseline và ASR+JSAR trên hai scale set.}
  \label{fig:final-term-metrics-by-scale-set}
\end{figure}

\begin{table}[H]
\centering
\caption{Kết quả định lượng giữa baseline và ASR+JSAR trên hai scale set}
\label{tab:final-term-scrfd-metrics}
\begingroup
\renewcommand{\arraystretch}{1.15}
\small
\begin{tabular}{lllll}
\toprule
\textbf{Scale set} & \textbf{Mô hình} & \textbf{Easy AP} & \textbf{Medium AP} & \textbf{Hard AP / mAP} \\
\midrule
Default source & Baseline & 0.9140 & 0.8969 & 0.7249 / 0.8453 \\
Default source & ASR+JSAR & 0.9039 & 0.8925 & 0.7690 / 0.8551 \\
Paper SR12 & Baseline & 0.9161 & 0.8999 & 0.7371 / 0.8510 \\
Paper SR12 & ASR+JSAR & 0.9028 & 0.8901 & 0.7738 / 0.8556 \\
\bottomrule
\end{tabular}
\endgroup
\end{table}

Delta của ASR+JSAR so với baseline được minh hoạ ở Hình~\ref{fig:final-term-asr-delta}.

\begin{figure}[H]
  \centering
  \includegraphics[width=0.88\textwidth]{final_term_spos_scrfd/asr_delta_by_scale_set.png}
  \caption{Delta của ASR+JSAR so với baseline trên từng chỉ số. Lợi ích lớn nhất luôn tập trung ở Hard AP, trong khi Easy AP và Medium AP giảm nhẹ.}
  \label{fig:final-term-asr-delta}
\end{figure}

\section{Vì sao ASR+JSAR tốt hơn baseline ở cả hai scale set}

\subsection{JSAR tăng mật độ supervision cho tiny faces một cách ổn định}

Ở cả hai scale set, JSAR đều làm tăng rõ rệt số positives trên mỗi tiny GT, như Hình~\ref{fig:final-term-jsar-tiny-supervision}.

\begin{figure}[H]
  \centering
  \includegraphics[width=0.88\textwidth]{final_term_spos_scrfd/jsar_tiny_supervision.png}
  \caption{Mật độ positive assignment trên tiny GT trước và sau JSAR. Cải tiến này giữ hiệu lực ổn định ở cả default scale set lẫn paper SR12.}
  \label{fig:final-term-jsar-tiny-supervision}
\end{figure}

Các trị số chính là:
\begin{itemize}
  \item \textbf{Default source}: tiny positives / GT tăng từ \(1.73\) lên \(2.67\), tương ứng boost ratio \(1.542\times\).
  \item \textbf{Paper SR12}: tiny positives / GT tăng từ \(1.76\) lên \(2.67\), tương ứng boost ratio \(1.520\times\).
\end{itemize}

Điều này cho thấy lợi ích của JSAR không đến từ nhiễu ngẫu nhiên trong huấn luyện.
Nó đến từ việc tiny faces thực sự nhận được supervision dày hơn ở giai đoạn assignment.

\subsection{Gain trên hard subset tập trung đúng vào tiny faces}

Hình~\ref{fig:final-term-hard-size-compare} cho thấy recall gain trên hard subset luôn dồn mạnh vào hai bin nhỏ nhất:
\([0,8)\) và \([8,16)\).

\begin{figure}[H]
  \centering
  \includegraphics[width=0.88\textwidth]{final_term_spos_scrfd/hard_size_gain_by_scale_set.png}
  \caption{Recall gain theo hard-size bin của ASR+JSAR so với baseline. Ở cả hai scale set, phần lợi ích gần như dồn vào hai vùng mặt rất nhỏ.}
  \label{fig:final-term-hard-size-compare}
\end{figure}

Chi tiết cho từng scale set được thể hiện ở Hình~\ref{fig:final-term-hard-size-solo}.

\begin{figure}[H]
  \centering
  \begin{subfigure}[t]{0.48\textwidth}
    \centering
    \includegraphics[width=\textwidth]{final_term_spos_scrfd/default_hard_size_gain.png}
    \caption{Default source scale set}
  \end{subfigure}
  \hfill
  \begin{subfigure}[t]{0.48\textwidth}
    \centering
    \includegraphics[width=\textwidth]{final_term_spos_scrfd/paper_hard_size_gain.png}
    \caption{Paper SR12 scale set}
  \end{subfigure}
  \caption{Recall gain theo hard-size bin cho từng scale set.}
  \label{fig:final-term-hard-size-solo}
\end{figure}

Với default scale set:
\begin{itemize}
  \item \([0,8)\): \(+0.138\)
  \item \([8,16)\): \(+0.119\)
  \item \([16,32)\): \(+0.009\)
  \item từ \(32\) px trở lên: gần \(0\) hoặc âm nhẹ.
\end{itemize}

Với paper SR12:
\begin{itemize}
  \item \([0,8)\): \(+0.108\)
  \item \([8,16)\): \(+0.104\)
  \item \([16,32)\): \(+0.003\)
  \item từ \(32\) px trở lên: âm nhẹ.
\end{itemize}

Như vậy, lợi ích của ASR+JSAR là lợi ích rất mục tiêu: nó tập trung vào đúng miền dữ liệu mà SCRFD khó nhất, tức là hard tiny faces.

\subsection{ASR và JSAR bổ sung cho nhau}

Từ cơ chế và kết quả trên, có thể nhìn ASR và JSAR như hai thành phần giải quyết hai nút thắt khác nhau:
\begin{itemize}
  \item \textbf{ASR} đổi phân phối sample mà detector thường xuyên nhìn thấy trong quá trình train.
  \item \textbf{JSAR} đổi mật độ target supervision mà tiny faces nhận được tại bước assignment.
\end{itemize}

Nếu chỉ có ASR, tiny faces có thể vẫn bị under-assigned.
Nếu chỉ có JSAR, detector chưa chắc đã nhìn thấy đủ nhiều tình huống tiny/hard để khai thác hết lợi ích.
Hai thành phần chỉ phát huy rõ khi được dùng cùng nhau.

\section{Cách hiểu khác biệt giữa default scale set và paper SR12}

\subsection{Paper SR12 có static scale pool tiny-oriented hơn}

Hình~\ref{fig:final-term-scale-pool} cho thấy hình học của hai scale pool.

\begin{figure}[H]
  \centering
  \includegraphics[width=0.92\textwidth]{final_term_spos_scrfd/scale_pool_geometry.png}
  \caption{Hình học của hai scale pool và phân phối xác suất scale. Paper SR12 có xu hướng ít zoom-in mạnh hơn nhưng lại dành nhiều khối lượng hơn cho vùng zoom-out, do đó baseline đã tiny-oriented hơn.}
  \label{fig:final-term-scale-pool}
\end{figure}

Một số trị số quan trọng:
\begin{itemize}
  \item \textbf{Default baseline}: \(E[1/\text{scale}] = 1.270\), khối lượng extreme zoom-in \((\text{scale}\le 0.6)\) là \(0.300\), extreme zoom-out \((\text{scale}\ge 2.0)\) là \(0.100\).
  \item \textbf{Paper SR12 baseline}: \(E[1/\text{scale}] = 0.890\), khối lượng extreme zoom-in là \(0.083\), extreme zoom-out là \(0.250\).
\end{itemize}

Điều đó cho thấy static scale pool SR12 vốn đã ít zoom-in mạnh hơn nhưng nhiều zoom-out mạnh hơn, nên tự nhiên đẩy train distribution về phía nhiều tiny faces hơn ngay từ baseline.

\subsection{Phân phối mặt sau sample redistribution của paper SR12 đã rất tiny-heavy}

Hình~\ref{fig:final-term-distribution-pressure} cho thấy paper SR12 baseline đã nhìn thấy nhiều tiny faces hơn đáng kể so với default baseline.

\begin{figure}[H]
  \centering
  \includegraphics[width=0.92\textwidth]{final_term_spos_scrfd/distribution_pressure.png}
  \caption{So sánh áp lực phân phối sau sample redistribution. Paper SR12 baseline đã tiny-oriented hơn rõ rệt, nên khoảng trống để ASR cải thiện thêm trở nên nhỏ hơn.}
  \label{fig:final-term-distribution-pressure}
\end{figure}

So với default scale set:
\begin{itemize}
  \item \textbf{Default baseline}: tiny ratio \(=0.475\), median face size \(=17.12\), tiny \(\rightarrow \ge 16\) px promotion ratio \(=0.223\).
  \item \textbf{Paper SR12 baseline}: tiny ratio \(=0.582\), median face size \(=12.97\), tiny \(\rightarrow \ge 16\) px promotion ratio \(=0.092\).
\end{itemize}

Nói cách khác, baseline của paper SR12 vốn đã nhìn thấy nhiều tiny faces hơn, nhưng đồng thời cũng ít ``promote'' tiny faces sang vùng dễ học hơn.
Đó là lý do vì sao:
\begin{itemize}
  \item paper SR12 baseline mạnh hơn default baseline trên Hard;
  \item nhưng gain biên của ASR+JSAR so với baseline lại nhỏ hơn một chút.
\end{itemize}

\subsection{Kết luận giữa hai scale set}

Các kết quả hiện tại không ủng hộ nhận định rằng paper SR12 tệ hơn.
Ngược lại:
\begin{itemize}
  \item \textbf{paper SR12 baseline} mạnh hơn \textbf{default baseline};
  \item \textbf{paper SR12 + ASR+JSAR} gần như ngang và hơi nhỉnh hơn \textbf{default ASR+JSAR};
  \item phần chênh lệch chính nằm ở chỗ scale pool SR12 đã tiny-oriented hơn ngay từ đầu.
\end{itemize}

Do đó, khi static scale pool đã mang sẵn thiên hướng tiny-focused, baseline trở nên mạnh hơn và \emph{headroom} để ASR+JSAR cải thiện thêm sẽ nhỏ hơn.
Tuy nhiên cơ chế gain của phương pháp vẫn giữ nguyên: lợi ích vẫn tập trung vào tiny/hard faces và đến từ việc đổi cả phân phối sample lẫn mật độ assignment.

\section{Kết luận}

Phần thực nghiệm bổ sung này đưa ra ba kết luận chính:
\begin{itemize}
  \item \textbf{ASR+JSAR} tốt hơn baseline ở cả hai scale set vì tăng training signal đúng vào tiny/hard faces.
  \item \textbf{JSAR} là thành phần trực tiếp làm tăng mật độ supervision cho tiny faces, còn \textbf{ASR} là thành phần dịch phân phối sample về phía nhiều tình huống tiny/hard hữu ích hơn.
  \item \textbf{Paper SR12} không làm mô hình tệ đi; ngược lại, nó tạo ra baseline mạnh hơn do static scale pool tiny-oriented hơn, đồng thời làm phần gain biên của ASR+JSAR nhỏ lại vì baseline đã mang sẵn một phần lợi ích đó.
\end{itemize}

Nói ngắn gọn, cải tiến ASR+JSAR là ổn định theo đúng mục tiêu thiết kế: ưu tiên tiny/hard faces.
Hiệu quả tuyệt đối của phương pháp vẫn phụ thuộc vào static scale pool ban đầu, nhưng cơ chế cải thiện cốt lõi không thay đổi.
